\subsection{Vaizdų paruošimas}
\label{subsec:vaizdu_paruosimas}

Pradinė „Muffin vs Chihuahua“ aibė atsisiunčiama jau padalinta į \texttt{train} ir \texttt{test} aplankus, kurių kiekvienas turi po du poaplankius pagal klases (\textit{chihuahua}, \textit{muffin}). Vaizdai yra skirtingų matmenų ir proporcijų, dalis – ne RGB. Kad juos būtų galima paduoti į konvoliucinį tinklą, klasė \texttt{ImageDataReader} taiko vienodą \texttt{torchvision.transforms} grandinę kiekvienam paveikslėliui:
\begin{itemize}
    \item \texttt{Resize($64 \times 64$)} – visi vaizdai pakeičiami į vienodą $64 \times 64$ raiškos formatą (konstanta \texttt{IMAGE\_SIZE} $= 64$).
    \item \texttt{ToTensor()} – $\texttt{PIL.Image}$ konvertuojama į \texttt{torch.Tensor} formatą, pikselių reikšmės iš $[0,\,255]$ perkeliamos į $[0,\,1]$ intervalą.
    \item \texttt{Normalize(mean, std)} – kanalų reikšmės standartizuojamos pagal \textit{ImageNet} statistikas: vidurkis $(0{,}485;\ 0{,}456;\ 0{,}406)$, standartinis nuokrypis $(0{,}229;\ 0{,}224;\ 0{,}225)$.
\end{itemize}

Validavimo aibė atskiriama klasėje \texttt{ImageDataSplitter}: iš kiekvienos klasės \texttt{train} aplanko atsitiktinai paimama $20\,\%$ failų ir perkeliama į naują \texttt{validation} aplanką, taip užtikrinant stratifikaciją pagal klases. Testavimo aibė lieka nepakitusi nuo pradinio šaltinio.

Skaičiavimo resursams taupyti buvo naudojama tik dalis duomenų iš visos duomenų aibės. Kiekvienos aibės dydis dauginamas iš parametro \texttt{subset\_ratio}, atrenkant atsitiktinius indeksus per \texttt{torch.randperm} su fiksuotu \textit{seed}. Hiperparametrų tyrimuose naudotas $\texttt{subset\_ratio} = 0{,}2$, o galutiniame geriausio modelio įvertinime – $\texttt{subset\_ratio} = 0{,}5$.
